\documentclass[a4paper]{article}

\usepackage[a4paper, top=8mm, bottom=8mm, left=8mm, right=8mm]{geometry}

\usepackage{polyglossia}
\setdefaultlanguage[babelshorthands=true]{russian}
\setotherlanguage{english}

\usepackage{fontspec}
\setmainfont{FreeSerif}
\newfontfamily{\russianfonttt}[Scale=0.7]{DejaVuSansMono}

\usepackage[tiny, compact]{titlesec}
\usepackage{titling}
\usepackage{hyperref}
\usepackage{bookmark}
\usepackage{csquotes}
\usepackage{amsmath}
\usepackage{booktabs}
\usepackage{array}

\setlength{\parindent}{1.25em}
\setlength{\parskip}{0pt}

\titleformat{\section}
{\normalfont\bfseries\large}
{\thesection.}
{0.5em}
{}

\titleformat{\subsection}
{\normalfont\bfseries\normalsize}
{\thesubsection.}
{0.5em}
{}

\hypersetup{
pdftitle={Исследование способов оптимизации операции перемножения плотных матриц (SGEMM) на архитектуре AMD RDNA3 с использованием OpenCL},
pdfauthor={Гарин Александр Евгеньевич},
hidelinks
}

\begin{document}

\begin{center}
{\Large\bfseries Исследование способов оптимизации операции перемножения плотных матриц (SGEMM) на архитектуре AMD RDNA3 с использованием OpenCL}

```
\vspace{0.5em}

Гарин Александр Евгеньевич

Группа 25.Б71-мм

Кафедра системного программирования

Руководитель: Григорьев Семен Вячеславович

Семестр: 2
```

\end{center}

\section{Постановка задачи}

Работа посвящена исследованию способов оптимизации операции SGEMM (Single-precision General Matrix Multiplication) --- перемножения плотных матриц с числами одинарной точности --- на интегрированном графическом процессоре AMD Radeon 760M с архитектурой RDNA3 с использованием OpenCL (Open Computing Language) --- открытого стандарта для программирования гетерогенных вычислительных систем.

В качестве основы для реализации использовался учебный материал Cedric Nugteren, посвящённый оптимизации SGEMM с использованием OpenCL. Исходный материал содержит последовательность реализаций операции матричного умножения, в которых последовательно применяются различные способы оптимизации.

\footnote{URL: \url{https://cnugteren.github.io/tutorial/pages/page1.html} (дата обращения: 05.09.2026).}

Целью работы является исследование производительности различных вариантов реализации SGEMM на архитектуре AMD RDNA3 и подбор параметров для одного из вариантов реализации.

Для достижения цели были поставлены следующие задачи:
\begin{itemize}
\item реализовать и запустить выбранные варианты SGEMM;
\item провести измерения времени выполнения и производительности;
\item сравнить результаты различных вариантов реализации;
\item исследовать влияние параметров блокирования на производительность выбранной реализации;
\item определить конфигурацию параметров, показавшую наибольшую производительность среди исследованных вариантов.
\end{itemize}

\section{Описание предлагаемого решения}

Для проведения исследования были выбраны реализации SGEMM с номерами 1, 2, 3, 4, 5, 6, 7 и 10 из учебного материала. Последовательность реализаций позволяет сравнить несколько вариантов организации вычислений и доступа к данным.

В первой реализации вычисления выполняются без применения рассматриваемых оптимизаций. В последующих реализациях используются различные способы организации доступа к данным, локальная память и блокирование вычислений.

Для исследования параметров реализации 4 использовался скрипт \texttt{extra/tuner.sh}, предназначенный для перебора значений параметров \texttt{TS} и \texttt{WPT}. Параметр \texttt{TS} определяет размер блока матричного умножения, а \texttt{WPT} (Work Per Thread) --- количество элементов результата, вычисляемых одним потоком.

Исходный учебный материал использует для соответствующей реализации значения \texttt{TS=32} и \texttt{WPT=8}. В рамках исследования дополнительно проверялись значения, уменьшенные в два раза: \texttt{TS=16} и \texttt{WPT=4}.

Некоторые варианты реализации из исходного учебного материала при запуске приводили к зависанию системы или её перезагрузке. Причина такого поведения в рамках данной работы отдельно не исследовалась, поэтому причинно-следственная связь между конкретной реализацией и наблюдавшимся поведением не устанавливается.

Исходный код исследования размещён в репозитории проекта.

\footnote{URL: \url{https://github.com/aleksandrgarin/myGEMM/tree/practice-optimization} (дата обращения: 05.09.2026).}

\section{Эксперименты}

Эксперименты проводились на ноутбуке Honor MagicBook X16 (BORN-H5651) с процессором AMD Ryzen 5 7640HS и интегрированным графическим процессором AMD Radeon 760M.

Аппаратная и программная конфигурация имела следующие характеристики:
\begin{itemize}
\item архитектура графического процессора (GPU, Graphics Processing Unit) --- RDNA3;
\item идентификатор GPU --- \texttt{gfx1103};
\item операционная система --- Ubuntu 24.04.4 LTS;
\item компилятор GCC/G++ --- версия 13.3.0;
\item ROCm (Radeon Open Compute) --- версия 6.4.4;
\item среда выполнения OpenCL --- \texttt{rocm-opencl-runtime} версии 6.4.4;
\item платформа OpenCL --- AMD Accelerated Parallel Processing, OpenCL 2.1 AMD-APP (3649.0);
\item поддерживаемая устройством версия OpenCL --- 2.0;
\item версия драйвера и среды выполнения --- 3649.0 (HSA1.1, LC), где HSA (Heterogeneous System Architecture) --- архитектура взаимодействия центрального и графического процессоров, а LC --- режим компиляции вычислительных программ;
\item загрузчик OpenCL ICD (Installable Client Driver) --- версия 2.3.2;
\item профиль OpenCL --- 3.0;
\item ядро Linux --- версия 6.17.0-1020-oem.
\end{itemize}

Запуски вычислений выполнялись в режиме максимальной производительности системы.

Компиляция исходного кода выполнялась с использованием стандарта C++11 и параметров \texttt{-O3} и \texttt{-Wall}.

\subsection{Входные данные}

Для экспериментов использовались три квадратные плотные матрицы размера $4096 \times 4096$ с элементами типа \texttt{float}. Такой размер матриц использовался в исходном учебном материале, на основе которого построена исследуемая реализация.

Элементы матриц инициализировались по формулам, использованным в исходном материале:

$$
A_i = 3.6i + i^2 + 3.1,
$$

$$
B_i = 1.2i + 0.01i^2 + 13.9.
$$

Матрица $C$ перед выполнением вычислений заполнялась нулевыми значениями.

Каждая реализация запускалась 10 последовательных раз. Для оценки производительности использовалось среднее арифметическое времени выполнения. Средняя производительность также вычислялась как среднее арифметическое результатов отдельных запусков. Запуски, показавшие отдельные значения времени или производительности, из расчёта не исключались.

В качестве дополнительной характеристики разброса результатов использовалась половина размаха:

$$
\Delta G = \frac{G_{\max} - G_{\min}}{2},
$$

где $G_{\max}$ и $G_{\min}$ --- максимальное и минимальное значения производительности среди десяти запусков.

Производительность измерялась в GFLOPS (Giga Floating Point Operations Per Second) и рассчитывалась по формуле:

$$
G = \frac{2MNK}{t \cdot 10^9},
$$

где $M$, $N$ и $K$ --- размеры матриц, а $t$ --- среднее время выполнения в секундах. Множитель 2 учитывает операции умножения и сложения, выполняемые для каждого элемента результата.

\subsection{Сравнение реализаций}

Результаты измерений времени выполнения и производительности представлены в таблице.

\begin{center}
\begin{tabular}{c >{\raggedright\arraybackslash}p{7.5cm} r r}
\toprule
№ & Реализация & Время, с & Производительность, GFLOPS \
\midrule
1  & Базовая реализация & 3.626 & $37.9 \pm 0.2$ \
2  & Блокирование вычислений & 0.521 & $263.4 \pm 14.5$ \
3  & Использование локальной памяти & 0.303 & $452.7 \pm 29.4$ \
4  & Блокирование и оптимизация доступа к данным & 0.297 & $461.2 \pm 30.1$ \
5  & Изменённая организация вычислений & 0.460 & $298.4 \pm 15.1$ \
6  & Одномерная обработка нескольких элементов & 0.323 & $425.4 \pm 36.2$ \
7  & Дополнительная оптимизация доступа к данным & 0.310 & $443.2 \pm 38.0$ \
10 & Оптимизированная организация вычислений & 0.424 & $316.0 \pm 15.0$ \
\bottomrule
\end{tabular}
\end{center}

Наибольшую производительность среди исследованных реализаций показала реализация 4 --- $461.2$~GFLOPS. Для сравнения, базовая реализация 1 показала $37.9$~GFLOPS. Таким образом, производительность реализации 4 выше производительности реализации 1 примерно в $12.2$~раза.

Реализация 5 показала производительность $298.4$~GFLOPS, что на $35.3$~% ниже результата реализации 4. Это показывает, что различия в организации вычислений и доступа к данным могут существенно влиять на получаемую производительность, однако вклад отдельных применённых оптимизаций в рамках данного эксперимента отдельно не измерялся.

\subsection{Настройка параметров реализации 4}

Для реализации 4 дополнительно исследовалось влияние двух параметров: \texttt{TS}, задающего размер блока, и \texttt{WPT}, задающего количество элементов результата, вычисляемых одним потоком.

В исходном учебном материале для рассматриваемой реализации используются значения \texttt{TS=32} и \texttt{WPT=8}. Для исследования были выбраны значения

$$
TS \in \{16,32\},
\qquad
WPT \in \{4,8\}.
$$

Таким образом, было проверено четыре комбинации параметров.

\begin{center}
\begin{tabular}{c c r}
\toprule
$TS$ & $WPT$ & Производительность, GFLOPS \
\midrule
16 & 4 & $462.7 \pm 50.8$ \
16 & 8 & $474.0 \pm 31.2$ \
32 & 4 & $477.0 \pm 45.1$ \
32 & 8 & $433.9 \pm 37.2$ \
\bottomrule
\end{tabular}
\end{center}

Наибольшую производительность среди проверенных комбинаций показала конфигурация $TS=32$, $WPT=4$ --- $477.0$~GFLOPS.

Исследование не являлось полным перебором возможных значений параметров: проверялись только значения $TS \in {16,32}$ и $WPT \in {4,8}$. Поэтому полученная конфигурация является лучшей только среди исследованных комбинаций.

\section{Ограничения проведённого исследования}

В работе исследовались только выбранные реализации SGEMM и ограниченный набор значений параметров \texttt{TS} и \texttt{WPT}. Полный перебор возможных параметров не проводился.

Размер матриц $4096 \times 4096$ и количество последовательных запусков, равное 10, были выбраны в соответствии с исходным учебным материалом, на основе которого построена работа. Поэтому полученные результаты характеризуют производительность исследуемых реализаций при заданном размере входных данных и выбранных условиях эксперимента.

В рамках исследования не проводилось отдельное измерение вклада каждой оптимизации в итоговое ускорение. Сравнение проводилось между готовыми вариантами реализаций, последовательно представленными в исходном учебном материале.

\section{Заключение}

В ходе работы были исследованы варианты реализации SGEMM с использованием OpenCL на интегрированном графическом процессоре AMD Radeon 760M с архитектурой RDNA3.

Для выбранных реализаций были проведены измерения времени выполнения и производительности на матрицах размера $4096 \times 4096$. Наибольшую производительность среди основных исследованных вариантов показала реализация 4 --- $461.2$~GFLOPS, что примерно в $12.2$~раза выше результата базовой реализации 1, показавшей $37.9$~GFLOPS.

Для реализации 4 был проведён дополнительный перебор значений параметров $TS$ и $WPT$. Среди четырёх проверенных комбинаций наибольшую производительность показала конфигурация $TS=32$, $WPT=4$ --- $477.0$~GFLOPS.

При этом исследование не охватывало полный диапазон возможных параметров и не позволяло определить вклад отдельных оптимизаций в итоговую производительность.

\end{document}
